\begin{algorithm}[H]
\caption{\GradeDP: cumulative-grade interval feasibility or weighted delay}
\label{alg:grade-dp}
\KwIn{Ordinary intervals; $r$ capacities $b_0<\cdots<b_{r-1}$;
feasibility deadlines, or weights, delay levels $d_j=b_j-1$,
and zero-pattern mass $p_0$.}
\KwOut{A feasible graded selection, or a minimum surrogate-delay selection.}
Initialize the state $(c_0,\ldots,c_{r-1};s_0,\ldots,s_{r-1})$ to all zeros,
with cost zero; all other states are unreachable\;
\For{$q=1,\ldots,M$}{
  From each reachable state, try skipping $q$ or selecting it at grade $i$\;
  A grade-$i$ selection increments $c_j$ and sets $s_j=q$ for all $j\ge i$\;
  Reject any transition with $c_j>b_j$\;
  \ForEach{interval $I=[l,q]$ ending here}{
    \eIf{mode is feasibility}{
      Reject if $s_j<l$ for its deadline grade $j$\;
    }{
      Find the least $j$ with $s_j\ge l$; reject if none\;
      Add $2w(I)d_j$ to this transition's cost\;
    }
  }
  Keep one feasible witness, or the minimum cost, for each resulting state\;
}
In weighted mode add $2p_0c_{r-1}$ to each final state's cost\;
Reconstruct a winning selection and order it by grade\;
\KwRet the selection and schedule, or infeasible\;
\end{algorithm}

For circular demands enumerate the forced first candidate, remove
its demands, cut there, and reduce \emph{every} capacity by one.
Residual expected costs still use the original delay levels and add
$2p_0(1+c_{r-1})$ at completion. In maximum-depth mode append a final
grade with the total-count cap if it is not already a capacity.
Section~\ref{sec:quasipolynomial} proves the state invariant and
the $(M+1)^{O(r)}(B+1)^{O(r)}$ bound.
